\documentclass[11pt]{article}
\usepackage[margin=1in]{geometry}
\usepackage{setspace}
\usepackage{enumitem}
\setlist{noitemsep, topsep=0pt}

\begin{document}

\begin{center}
\textbf{\large INFORMED CONSENT TO PARTICIPATE IN RESEARCH}\\[0.5em]
\textbf{From "What If" to "Now What": Improving Requirement Engineering Among End-User Programmers Through Exploratory Testing}
\end{center}

\noindent\textbf{Principal Investigator:} Hudson Mitchell-Pullman, UC San Diego\\
\textbf{Contact:} hudsonmitchellpullman@gmail.com | (858) 868-0309

\vspace{1em}

\noindent\textbf{Purpose}\\
This study seeks to understand how novice computer science students approach requirements engineering through the testing of their software systems. The research compares outcomes between students who use a tutoring system designed to enhance high-level testing skills and a control group. A qualitative rubric will be used to evaluate whether there are meaningful differences in requirements engineering outcomes between the two groups.

\vspace{0.8em}

\noindent\textbf{Procedures}\\
If you agree to participate, you will complete a 20-minute pre-assessment, a 40-minute learning activity with a large language model (LLM), and a 20-minute post-assessment. After these activities, you will participate in a 10-minute one-on-one interview.

\vspace{0.8em}

\noindent\textbf{Risks}\\
The risks associated with this study are minimal and similar to those found in typical learning environments. The learning intervention and the control group may have different levels of effectiveness, which could result in varying learning outcomes. To address this, all code will be made open-source after the study, and participants in the control group will be granted free access to the tutoring system, with API credits subsidized by the research team.

Additionally, due to the unpredictable nature of large language models (LLMs), there is a potential risk that the LLM may generate unwanted or inappropriate responses that could negatively affect your learning experience. The prompts used in this study have been rigorously tested to minimize this risk. Participation in this in-person study will require approximately 90 minutes of your time.

\vspace{0.8em}

\noindent\textbf{Benefits}\\
This study is expected to enhance exploratory testing and requirements engineering skills, which are relevant to both end-user and AI-assisted programming. The learning intervention is grounded in empirical evidence from cognitive science, and all necessary measures have been taken to ensure its pedagogical validity.

\vspace{0.8em}

\noindent\textbf{Confidentiality}\\
Your responses will be kept anonymous, and any identifiable information will be redacted. Data will be retained indefinitely, but will not be shared with any non-IRB-approved entity.

\vspace{0.8em}

\noindent\textbf{Compensation}\\
All participants will receive a $\$15$ Amazon gift card, which will be distributed via email within one week after the study concludes. Participants who refer a friend who enrolls in the study will receive an additional $\$5$ bonus, regardless of the number of referrals.

\vspace{0.8em}

\noindent\textbf{Voluntary Participation}\\
Participation is voluntary. You may refuse to participate or withdraw at any time without penalty.

\vspace{0.8em}

\noindent\textbf{Questions}\\
If you have questions about this study, contact Hudson Mitchell-Pullman at hudsonmitchellpullman@gmail.com. For questions about your rights as a research participant, contact GSDSEF SRC.

\vspace{1.5em}

\noindent\textbf{Consent}\\
By signing below, I confirm that I have read this form, had the opportunity to ask questions, and voluntarily agree to participate.

\vspace{1.5em}

\noindent\begin{tabular}{@{}p{0.32\textwidth}p{0.02\textwidth}p{0.32\textwidth}p{0.02\textwidth}p{0.32\textwidth}@{}}
\hrulefill & & \hrulefill & & \hrulefill \\
Participant Signature & & Printed Name & & Date
\end{tabular}

\end{document}
